\chapter{Methodology}
\label{ch:methodology}

This chapter describes the data, the exploratory analysis that shaped the design, the eight
model pipelines, the protocol used to train and evaluate them, and the system that serves
the result.

\section{Dataset}
\label{sec:dataset}

The project uses the public \gls{ulb} credit card dataset, which contains transactions made
by European cardholders over two days in September 2013. Each record has a \texttt{Time}
value, an \texttt{Amount} value, twenty-eight anonymised PCA components \texttt{V1}--
\texttt{V28}, and a binary \texttt{Class} label.

Before anything else, 1{,}081 exact duplicate rows are removed, leaving 283{,}726
transactions of which 473 are fraudulent. De-duplication is done before splitting on
purpose, because identical rows landing in both the training and test sets would quietly
inflate every reported score. Table~\ref{tab:profile} gives the full profile.

\begin{table}[htbp]
    \centering
    \caption{Dataset profile.}
    \label{tab:profile}
    \small
    \begin{tabular}{l r}
\hline
\textbf{Statistic} & \textbf{Value} \\
\hline
Transactions after de-duplication & 283{,}726 \\
Exact duplicate rows removed & 1{,}081 \\
Fraudulent transactions & 473 \\
Legitimate transactions & 283{,}253 \\
Fraud rate (\%) & 0.167 \\
Imbalance ratio (legit : fraud) & 598.844 \\
Missing values & 0 \\
Features used for modelling & 29 \\
Median amount, legitimate (EUR) & 22 \\
Median amount, fraud (EUR) & 9.820 \\
Mean amount, legitimate (EUR) & 88.414 \\
Mean amount, fraud (EUR) & 123.872 \\
95th percentile amount, legitimate (EUR) & 365 \\
95th percentile amount, fraud (EUR) & 655.820 \\
Largest fraudulent amount (EUR) & 2125.870 \\
Hour of day with highest fraud rate & 2 \\
Peak hourly fraud rate (\%) & 1.451 \\
Lowest hourly fraud rate (\%) & 0.048 \\
Strongest single feature & V17 \\
Its absolute correlation with the label & 0.313 \\
Second strongest feature & V14 \\
Its absolute correlation with the label & 0.293 \\
\hline
\end{tabular}
\end{table}

From the table we can see that only 473 of the 283{,}726 remaining transactions are
fraudulent, that no values are missing anywhere in the data, and that 29 of the columns are
used for modelling. All of these values are computed by
\texttt{src/\allowbreak deep\_\allowbreak analysis.py} and written to
\texttt{reports/\allowbreak dataset\_\allowbreak profile.csv}.

The imbalance ratio of 598.8 legitimate transactions per fraud is the single most important
property of this dataset, and it drives every design decision that follows.
Figure~\ref{fig:balance} shows it on a logarithmic axis, because on a linear scale the
fraud bar is not visible at all.

\begin{figure}[htbp]
    \centering
    \includegraphics[width=0.62\textwidth]{01_class_balance.png}
    \caption{Class balance.}
    \label{fig:balance}
\end{figure}

From the figure above we can see how far apart the two bars are even after the axis has
been compressed logarithmically. On a linear scale the fraud bar would be about one pixel
high, and this gap is the core difficulty of the problem.

\section{Exploratory Data Analysis}
\label{sec:eda}

\subsection{Transaction Amounts}

A common intuition is that fraud chases large transactions. This dataset contradicts that.
The median fraudulent transaction is EUR 9.82 against EUR 22.00 for legitimate ones, so
fraud is typically smaller. The mean tells the opposite story (EUR 123.87 against EUR
88.41), because a small number of large frauds pull the mean upward while most fraudulent
activity sits at very low values. That is consistent with card testing, where an attacker
checks a stolen card with a trivial purchase before attempting anything larger.

Figure~\ref{fig:amount} shows this as a cumulative distribution, which compares the two
groups more clearly than overlaid histograms would: the curve sitting higher and further
left belongs to the smaller-valued group.

\begin{figure}[htbp]
    \centering
    \includegraphics[width=0.78\textwidth]{02_amount_by_class.png}
    \caption{Transaction amount by class.}
    \label{fig:amount}
\end{figure}

From the figure above we can see that the fraud curve lies above and to the left of the
legitimate curve across almost the whole range, which means fraudulent transactions are
typically the smaller ones. This rules out any simple amount-threshold rule.

\subsection{Temporal Pattern}

Grouping by hour of day shows a strong pattern. Legitimate volume collapses overnight while
fraudulent activity does not, so the fraud rate peaks at 1.451 percent around 02:00, about
nine times the overall rate of 0.167 percent. Its lowest point is 0.048 percent during the
busiest daytime hours, a spread of a factor of thirty. Figure~\ref{fig:hourly} plots both
quantities against the hour of day.

\begin{figure}[htbp]
    \centering
    \includegraphics[width=0.8\textwidth]{03_hourly_pattern.png}
    \caption{Hourly transaction pattern.}
    \label{fig:hourly}
\end{figure}

From the figure above we can see that the two panels move in opposite directions: the
overnight hours where volume is lowest in the top panel are the same hours where the fraud
rate peaks in the bottom panel. The panels are stacked rather than plotted on twin axes,
which would imply a relationship the data does not establish.

This is a real signal, and the model deliberately does not use it, for the reason given in
Section~\ref{sec:features}: \texttt{Time} in this dataset cannot be mapped to a real clock
for an unseen transaction. A production system with true timestamps should revisit this
decision.

\subsection{Feature Signal}

Ranking each feature by the absolute Pearson correlation between its value and the label
puts V17 (0.313), V14 (0.293), V12 (0.251), and V10 (0.207) well clear of the rest. Because
the components are anonymised, this ranking is the only description available of which
inputs matter. Figure~\ref{fig:signal} shows the full ranking.

\begin{figure}[htbp]
    \centering
    \includegraphics[width=0.6\textwidth]{04_feature_signal.png}
    \caption{Feature signal ranking.}
    \label{fig:signal}
\end{figure}

From the figure above we can see that the signal is concentrated in a handful of
components: the top four stand well clear, and the bars then decay quickly towards zero.
Figure~\ref{fig:dists} looks at those four more closely.

\begin{figure}[htbp]
    \centering
    \includegraphics[width=0.88\textwidth]{05_top_feature_distributions.png}
    \caption{Top feature distributions.}
    \label{fig:dists}
\end{figure}

From the figure above we can see that fraud sits in a visibly different region of each of
the four distributions, and that separation is exactly what a classifier exploits.

\subsection{Correlation Structure and Separability}

Figure~\ref{fig:heatmap} shows how the features relate to one another.

\begin{figure}[htbp]
    \centering
    \includegraphics[width=0.82\textwidth]{06_correlation_heatmap.png}
    \caption{Feature correlation heatmap.}
    \label{fig:heatmap}
\end{figure}

From the figure above we can see that the off-diagonal is almost uniformly blank, so
V1--V28 are uncorrelated with each other, exactly as expected of PCA outputs. There is
therefore no redundant pair to drop and no case for further dimensionality reduction,
because the feature set is already orthogonal.

Finally, Figure~\ref{fig:cluster} plots all fraud cases against a 15{,}000-row sample of
legitimate ones on the two strongest features.

\begin{figure}[htbp]
    \centering
    \includegraphics[width=0.68\textwidth]{07_fraud_cluster.png}
    \caption{Fraud cluster.}
    \label{fig:cluster}
\end{figure}

From the figure above we can see that fraud (red) forms a loose but distinct cluster rather
than sitting inside the legitimate mass (blue). The two groups overlap, but they are
clearly distinguishable, which is direct visual evidence that a non-linear classifier has
real structure to learn here.

\section{Feature Selection}
\label{sec:features}

The model uses twenty-nine input features: the twenty-eight PCA components and
\texttt{Amount}. The \texttt{Time} field is left out on purpose. It is measured in seconds
elapsed since the first transaction in this two-day sample, so it means nothing for a
transaction observed later and cannot generalise. It is kept for exploratory analysis only,
where it is converted to hour of day.

No further feature engineering is applied. The PCA components are already decorrelated by
construction and their original meaning is hidden, so there is no basis for either
interaction terms or domain transformations.

\section{Pipeline Design}

Each of the eight models is a single pipeline built with scikit-learn
\citep{pedregosa2011scikit} and imbalanced-learn \citep{lemaitre2017imbalanced}, so the
saved winner is one self-contained object that the \gls{api} can call with a raw feature
table. Every pipeline chains together three steps in a fixed order: a scaler that
standardises the input features, an optional resampling step that corrects for class
imbalance, and a classifier that produces the fraud probability.

Packaging the three steps together is not just for tidiness. It is what guarantees that the
scaler fitted at training time is the exact scaler applied at inference, and that the
resampling step is never applied outside training.

The four algorithms are Logistic Regression, Random Forest, XGBoost, and LightGBM. Each is
trained twice, once with class weighting and once with SMOTE, giving eight pipelines. In
the weighted variants the classifier receives \texttt{class\_\allowbreak weight="balanced"} (or
\texttt{scale\_\allowbreak pos\_\allowbreak weight} for XGBoost) and no resampler is present. In the SMOTE variants
the classifier is left unweighted, because resampling has already balanced the classes and
applying both corrections would over-weight the minority class twice.

\section{Data Splitting Protocol}
\label{sec:splits}

The data is split once by stratified sampling into three parts, in the proportions 60/20/20.
Three splits are used rather than two because the default 0.5 probability cut-off is a poor
operating point at 0.167 percent positives, and choosing a better one is itself a decision
that can overfit. Table~\ref{tab:splits} states the sizes and the role of each partition.

\begin{table}[htbp]
    \centering
    \caption{Data splits.}
    \label{tab:splits}
    \small
    \begin{tabular}{l r r r l}
\hline
\textbf{Split} & \textbf{Rows} & \textbf{Frauds} & \textbf{Rate (\%)} & \textbf{Purpose} \\
\hline
Training & 170{,}235 & 284 & 0.167 & Fit model parameters \\
Validation & 56{,}745 & 94 & 0.166 & Choose the decision threshold \\
Test & 56{,}746 & 95 & 0.167 & Report final numbers (used once) \\
\hline
\end{tabular}
\end{table}

From the table we can see that stratification keeps the fraud rate the same in all three
partitions, and that the test set holds 95 of the 473 frauds. That set is used exactly
once, to produce the final reported numbers.

The protocol is therefore:
\begin{enumerate}
  \item Fit model parameters on the training set.
  \item Choose the decision threshold on the validation set.
  \item Report every final number on the test set, which took no part in either.
\end{enumerate}

\section{Threshold Selection}
\label{sec:threshold}

A classifier outputs a probability, and turning that into a decision requires a cut-off $t$.
The threshold is chosen to minimise the total number of mistakes on validation data:
\begin{equation}
  \label{eq:threshold}
  t^{\star} = \arg\min_{t} \Big( FN(t) + FP(t) \Big).
\end{equation}

Equation~\ref{eq:threshold} is written for this project rather than taken from a paper. It
is the equal-cost special case of the general cost-sensitive decision rule of
\citet{elkan2001foundations}, which selects the threshold minimising expected cost; setting
the cost of a missed fraud equal to the cost of a false alarm reduces that rule to counting
the two errors.

This objective was chosen over maximising $F_1$ because it is easier to explain: it
minimises the two counts a fraud team actually tracks, which are missed frauds and false
alarms. It is computed efficiently by sorting the validation probabilities once and
evaluating both counts through cumulative sums, giving $O(n \log n)$ overall.

There are two caveats. First, this objective assumes that one missed fraud costs exactly as
much as one false alarm, which is a modelling convenience rather than a business fact. A
monetary cost model is applied later to the same predictions to check how much difference
that assumption makes. Second, each model has a different probability distribution, so each
gets its own threshold. The tuned values range from 0.355 to 0.99999, which means comparing
raw probabilities across models is meaningless.

\section{Leakage Control}
\label{sec:leakage}

Three specific measures prevent optimistic bias:

\begin{enumerate}
  \item De-duplication before splitting. The 1{,}081 exact duplicates are removed
        first, so no row can appear in both training and test.
  \item Resampling inside the pipeline. SMOTE and the random samplers are pipeline
        steps, so \texttt{fit} applies them only to the training fold. Synthetic points
        derived from test frauds can therefore never enter training. This is the classic
        error that makes published results on this dataset impossible to reproduce.
  \item Threshold tuned on validation, not test. The test set influences nothing
        except the final report.
\end{enumerate}

As a safeguard, the training script reloads the saved winner from disk and re-scores the
test set, asserting that the PR-AUC matches to within $10^{-9}$. It also warns if the
PR-AUC exceeds 0.99, a value that on this dataset would indicate leakage rather than
success. The statistical analysis script repeats this check independently.

\section{Hyperparameter Search}

LightGBM's settings were selected by a randomised search over 40 candidate configurations,
scored by PR-AUC on the validation set. The pattern the search found was that many shallow
trees generalise better here than fewer deep ones: the selected configuration uses 1{,}500
estimators with only 15 leaves each. Table~\ref{tab:hyper} lists the final values.

\begin{table}[htbp]
    \centering
    \caption{LightGBM hyperparameters.}
    \label{tab:hyper}
    \small
    \begin{tabular}{l r l}
\hline
\textbf{Hyperparameter} & \textbf{Value} & \textbf{Role} \\
\hline
n\_estimators & 1500 & Number of boosting rounds \\
learning\_rate & 0.1 & Shrinkage applied to each tree \\
num\_leaves & 15 & Max leaves per tree (controls complexity) \\
min\_child\_samples & 30 & Min samples in a leaf \\
colsample\_bytree & 0.7 & Feature subsample per tree \\
reg\_lambda & 0.0 & L2 regularisation \\
class\_weight & balanced & Re-weights the loss by inverse class frequency \\
random\_state & 42 & Seed, fixed for reproducibility \\
\hline
\end{tabular}
\end{table}

From the table we can see the shape of the selected model: 1{,}500 small trees of at most
15 leaves each, a learning rate of 0.1 that keeps every individual tree's contribution
modest, and a 70 percent feature subsample per tree to limit overfitting.

The remaining algorithms use reasonable fixed settings rather than tuned ones (Random
Forest 200 trees; XGBoost 300 rounds at depth 5). This is a deliberate asymmetry and a
limitation: the comparison is between a tuned LightGBM and untuned competitors, which
biases it in LightGBM's favour.

\section{System Implementation}
\label{sec:system}

\subsection{Architecture}

The system splits into a modelling side written in Python and a serving side made up of an
HTTP API and a browser application.

% tabularx and its Y column come from projectreport.cls, so the description column wraps
% inside the text block instead of running off the right margin.
\begin{center}
\small
\begin{tabularx}{\textwidth}{@{}lY@{}}
\hline
\textbf{Stage} & \textbf{Responsibility} \\
\hline
\texttt{src/data.py}    & Load the CSV, de-duplicate, produce the stratified splits \\
\texttt{src/eda.py}     & Exploratory figures 1--7 \\
\texttt{src/models.py}  & Construct the eight pipelines \\
\texttt{src/train.py}   & Fit, tune thresholds, evaluate, persist the registry \\
\texttt{src/tune.py}    & Randomised hyperparameter search for LightGBM \\
\texttt{src/resampling\_experiment.py} & The controlled resampling study \\
\texttt{src/deep\_analysis.py} & The statistical analysis and its figures \\
\texttt{src/make\_tables.py} & Emit every LaTeX table in this report from the CSVs \\
\hline
\texttt{api/main.py}    & Load the registry once, serve scoring endpoints \\
\texttt{web/}           & React application \\
\hline
\end{tabularx}
\end{center}

The two sides communicate only through saved files, namely
\texttt{models/\allowbreak registry/\allowbreak *.joblib} and \texttt{models/\allowbreak metadata.json}. Models can therefore be
retrained without touching the application, and the application can be developed without
retraining models.

\subsection{Serving API}

The backend is a FastAPI service that loads all eight pipelines at startup. Every scoring
endpoint accepts an optional model name and threshold, so the caller chooses which trained
model scores the request and at what cut-off.

\begin{center}
\small
\begin{tabularx}{\textwidth}{@{}lY@{}}
\hline
\textbf{Endpoint} & \textbf{Purpose} \\
\hline
\texttt{GET /models}          & List all models with their test metrics \\
\texttt{GET /model-info}      & One model's features, threshold, metrics, and sweep \\
\texttt{GET /samples}         & Real held-out transactions for the interface to score \\
\texttt{POST /predict}        & Score a single transaction \\
\texttt{POST /predict/batch}  & Score an array of transactions \\
\texttt{POST /predict/csv}    & Score an uploaded file \\
\hline
\end{tabularx}
\end{center}

Rather than storing derived rates, \texttt{/model-info} returns each model's threshold sweep
as 201 raw confusion-count records of the form $\{t, TP, FP, FN, TN\}$ at
$t = 0.000, 0.005, \dots, 1.000$. Any metric, whether precision, recall, $F_1$, \gls{tpr} or
\gls{fpr}, is then derived in the browser without re-scoring anything. This keeps the API
small while giving the interface everything it needs.

\subsection{Web Application}

The frontend is a React application built with Vite and TypeScript. It has a landing page
and four working screens: score one transaction, score a batch, upload and score a
\gls{csv} file, and a dashboard.

A control bar applies to every screen, letting the user select any of the eight trained
models and set the decision threshold. Selecting a model resets the threshold to that
model's own tuned value, because Section~\ref{sec:threshold} established that thresholds
are not comparable across models. Since \texttt{V1}--\texttt{V28} are anonymised and cannot
be typed in by hand, the interface loads real held-out transactions for the user to score.
Figure~\ref{fig:detector} shows the single-transaction screen before anything is scored.

\begin{figure}[htbp]
    \centering
    \includegraphics[width=0.92\textwidth]{detector_page.png}
    \caption{Transaction scoring screen.}
    \label{fig:detector}
\end{figure}

From the figure above we can see the two controls that the report argues matter most. The
model selector names the pipeline in use together with its PR-AUC, and the threshold slider
sits at the tuned default of 0.636, labelled at each end with the trade-off it governs. The
four tabs across the top are the four screens, and the buttons above the feature form load a
real held-out fraud or legitimate transaction rather than asking the user to invent
anonymised values.

Scoring a transaction fills the empty panel on the right with a verdict.
Figure~\ref{fig:flagged} shows the result after a known fraud has been loaded and scored.

\begin{figure}[htbp]
    \centering
    \includegraphics[width=0.92\textwidth]{fraud_flagged.png}
    \caption{A flagged transaction.}
    \label{fig:flagged}
\end{figure}

From the figure above we can see that the model returns a fraud probability of 100.00
percent for this transaction and flags it, and that the badge confirms the verdict matches
the true label. The bar below the verdict places that probability on the 0 to 1 range with a
marker at the active threshold, so the user can see how far past the cut-off the score fell
rather than being given a bare yes or no.

\subsection{Three-Dimensional Visualisation}

Parts of this report describe objects that are naturally three-dimensional, and the
application renders them with WebGL using React Three Fiber.

The landing page is a single scroll-driven sequence in four stages, which walks a visitor
through the same path this report takes. It opens on a 3D credit card, then moves inside it
to a network of 31 nodes, one for each column of the dataset, where a node's distance from
the camera path is its permutation importance, so V14 and V4 pass closest. It then passes through a set of boosted decision trees, where each leaf
contributes a signed weight to a running total, and ends at a threshold gate that routes
scored transactions into an approved path and a flagged path. The transactions used in that
final stage are real held-out cases fetched from \texttt{GET /samples}, and a slider sets
the cut-off they are compared against, so moving it re-routes the stream on screen.
Figure~\ref{fig:landing} shows the opening stage of that sequence.

\begin{figure}[htbp]
    \centering
    \includegraphics[width=0.72\textwidth]{landing_page.png}
    \caption{Landing page.}
    \label{fig:landing}
\end{figure}

From the figure above we can see the first of the four stages, a 3D credit card rendered in
the browser with the prompt to scroll into its chip. The two buttons let a visitor either
begin the guided sequence or skip straight to the scoring screens, so the visualisation
never stands between a reviewer and the numbers.

The dashboard renders the decision surface, which plots the $8 \times 201$ grid of
threshold sweeps as a continuous landscape: threshold on one axis, model on the second,
$F_1$ as height, and precision as colour. The ridge running along the surface is each
model's usable operating band, and the cliff at low thresholds is where precision
collapses. A translucent plane tracks the threshold slider, so
moving the familiar 2D control walks a marker across the 3D surface. Beside it, the model
field renders the eight pipelines on their natural $4 \times 2$ grid of four algorithms by
two imbalance strategies, with bar height proportional to PR-AUC.

All of these views degrade gracefully. If WebGL is unavailable, or if the operating system
reports a preference for reduced motion, the application falls back to the two-dimensional
charts and tables, which remain the authoritative presentation of the numbers.

\section{Reproducibility}

Every random component is seeded with \texttt{random\_\allowbreak state = 42}, including the splitter,
the samplers, the bootstrap, the cross-validation folds, and each classifier. Package
versions are pinned in \texttt{requirements.txt} and listed in Appendix~\ref{app:repro}.
Re-running \texttt{src/train.py} followed by \texttt{src/\allowbreak deep\_\allowbreak analysis.py} and
\texttt{src/\allowbreak make\_\allowbreak tables.py} regenerates every number and figure in this report from the
raw CSV.
